\documentclass[12pt]{article}
\usepackage[utf8]{inputenc}
\usepackage[T1]{fontenc}
\usepackage{hyperref}
\title{Zaramo Language Resources for K-12 Curriculum Development}
\author{Compiled by Claude Sonnet 5.5 (claude-sonnet-5-5) for Andr\'as Kornai}
\date{}
\begin{document}
\maketitle

\section{Introduction}
Zaramo (ISO 639-3 zaj, Glottocode zara1247; also Kizaramo, Zalamo, Dzalamo) is a Bantu language of the Atlantic-Congo family, spoken mostly in villages around Dar es Salaam, Tanzania. About 1,000,000 speakers are estimated. The reference description says it is not known to be taught in schools. Wikipedia notes that very few translations exist apart from the 1975 New Testament.

\section{Primary Educational Resources}
The searches found no primers, textbooks or school materials. Only religious texts surfaced.

\section{Linguistic and Reference Materials}
\begin{itemize}
\item Wikipedia, ``Zaramo language''. \url{https://en.wikipedia.org/wiki/Zaramo_language}
\item Wikipedia, ``Zaramo people'': ethnographic background. \url{https://en.wikipedia.org/wiki/Zaramo_people}
\item Endangered Languages Project entry. \url{https://endangeredlanguages.com/node/88483}
\end{itemize}
No Zaramo dictionary or grammar turned up in these results.

\section{Religious and Community Resources}
\begin{itemize}
\item Zaramo Bible (2019), \emph{Lagano da Isambi da Mndewa na Mkombola Wetu Yesu Kilisto}, by Pioneer Bible Translators, on eBible.org in ePub, PDF and other formats. The licence is listed inconsistently across versions of the page, so check it before reuse. \url{https://ebible.org/Bible/details.php?id=zaj}
\item ScriptureEarth page listing a Bible app, a JESUS Film version and other scripture media. \url{https://www.scriptureearth.org/00eng.php?iso=zaj}
\end{itemize}

\section{Digital and Audio Resources}
Bible.is hosts a 1975 United Bible Societies recording, and ScriptureEarth lists an Android Bible app. These are the only audio or software items found. See the ScriptureEarth link above.

\section{Curriculum Development Context}
A very large speaker population lives in and around Dar es Salaam, but the published written material is limited to scripture. The 2019 Bible gives a modern written form of the language and a corpus for orthography work. Secular readers, a dictionary and a grammar would have to be developed. Swahili is the language of schooling in Tanzania.
\end{document}